\chapter*{Résumé}
\addcontentsline{toc}{chapter}{Résumé}

Ce projet tuteuré a pour objectif d'analyser le parallèle entre le traitement du langage par les modèles neuronaux de type Transformer et les mécanismes cognitifs du cerveau humain. Nous utilisons l'analyse de la dimensionnalité intrinsèque (ID) pour vérifier si l'IA effectue, comme le cerveau, une compression de l'information sémantique au cours de la lecture.

Pour répondre à cette problématique, nous formulons l'hypothèse que les phrases cohérentes créent une dynamique de «~ramping~» et un profil de compression spécifique (profil «~Hunchback~») qui seront absents ou modifiés lors du traitement de phrases «~Jabberwocky~» (syntaxiquement correctes mais dénuées de sens). Nous avons défini un protocole expérimental implémentant l'algorithme TwoNN sur le modèle GPT-2 pour mesurer ces variations géométriques couche par couche.

Dans le premier chapitre, nous explorons l'état de l'art en mettant en relation l'architecture des Transformers, le modèle neurolinguistique MUC et les propriétés mathématiques de l'ID. Puis, dans le chapitre suivant, nous détaillons notre méthodologie pour la phase d'expérimentation, en justifiant le choix de nos outils et la procédure de génération de nos stimuli linguistiques. Enfin, nous présentons dans le troisième chapitre notre feuille de route pour la mise en œuvre technique et l'analyse des résultats prévue au second semestre.